\section{Two-Element Array Theory and Design}
\label{sec:array-theory-design}

A \(2\times1\) patch array increases directivity by combining the fields from two patch elements.
If mutual coupling and element-pattern distortion are neglected, the total pattern can be approximated
by pattern multiplication:
\[
F_{\mathrm{total}}(\theta,\phi)
\approx
F_{\mathrm{element}}(\theta,\phi)AF(\theta,\phi),
\]
where \(F_{\mathrm{element}}\) is the single-patch element pattern and \(AF\) is the array factor.

For two identical elements along the \(x\)-axis with spacing \(d\), equal amplitudes, and progressive
phase shift \(\Delta\phi\), the normalized array factor in the array plane can be written as
\[
AF(\theta)
=
2\cos\left[\frac{1}{2}\left(k_0 d\sin\theta+\Delta\phi\right)\right],
\]
where \(k_0=2\pi/\lambda_0\). A broadside beam occurs for \(\Delta\phi=0\). Beam steering to
an angle \(\theta_0\) is obtained approximately from
\[
\boxed{\Delta\phi=-k_0 d\sin\theta_0.}
\]

The recommended starting center-to-center element spacing is
\[
d\approx\frac{\lambda_0}{2}\approx59.94~\mathrm{mm}
\]
for \(f_0=2.501~\mathrm{GHz}\). In practice, the spacing may need adjustment to reduce coupling,
avoid grating lobes, and fit the physical patch and feed geometry. The phase-shift values used as
beam-steering starting points are verified using the Python calculation in
Appendix~\ref{app:beamsteering-verification}, with the console output documented in
Listing~\ref{lst:beamsteering-console}.


For a two-element array with \(d=\lambda_0/2\), the ideal progressive phase shift required to steer the main beam toward \(\theta_0\) is
\[
\Delta\phi = -k_0 d \sin\theta_0 = -180^\circ \sin\theta_0 .
\]
The resulting initial phase-shift values are summarized in Table~\ref{tab:beamsteering-phase-starting-values}.

\begin{table}[H]
\centering
\caption{Initial progressive phase shifts for beam steering with \(d=\lambda_0/2\).}
\label{tab:beamsteering-phase-starting-values}
\small
\setlength{\tabcolsep}{8pt}
\renewcommand{\arraystretch}{1.20}

\begin{tabular}{@{}ccc@{}}
\toprule
\textbf{Target angle, \(\theta_0\)} &
\textbf{Phase-shift formula} &
\textbf{Progressive phase, \(\Delta\phi\)} \\
\midrule

\(0^\circ\) &
\(-180^\circ\sin(0^\circ)\) &
\(0^\circ\) \\

\(15^\circ\) &
\(-180^\circ\sin(15^\circ)\) &
\(-46.6^\circ\) \\

\(30^\circ\) &
\(-180^\circ\sin(30^\circ)\) &
\(-90.0^\circ\) \\

\(45^\circ\) &
\(-180^\circ\sin(45^\circ)\) &
\(-127.3^\circ\) \\

\bottomrule
\end{tabular}
\end{table}
